\section{中间步骤（Intermediate Steps）：推理的核心机制}

\subsection{历史脉络：从训练到提示}

Denny Zhou 梳理了中间步骤思想的发展历程：

\begin{enumerate}
    \item \textbf{2017 年，Ling et al.}：在论文中提出用自然语言推理步骤（rationale）求解数学题，训练 seq2seq 模型从头生成中间步骤再得出答案。被 Denny Zhou 称为"如同时间旅行者"。
    \item \textbf{2021 年，OpenAI GSM8K}：延续 2017 年思路，创建了包含中间步骤的数学数据集，用于 fine-tune GPT-3，大幅提升数学推理能力。
    \item \textbf{2021 年，Google Brain Scratchpad}：独立发现类似思路，但在程序合成（program synthesis）领域使用符号化中间步骤。
    \item \textbf{2022 年，Chain-of-Thought Prompting（Wei et al.）}：在 prompting 阶段广泛评估中间步骤的效果，展示了在几乎所有 NLP 任务上的显著提升。
\end{enumerate}

\begin{importantbox}{中间步骤是关键，而非具体方法}
无论是 training、fine-tuning 还是 prompting，真正起作用的是\textbf{中间步骤（intermediate steps）}本身。当提供包含中间步骤的示例时，LLM 会生成同样包含中间步骤的响应——这是 LLM 作为概率模型模仿输入模式的自然结果。
\end{importantbox}

\subsection{理论基础：为什么中间步骤有效}

Denny Zhou 引用了 2024 年与 Stanford 的 Branden Srouji 合作的理论工作：

\begin{itemize}
    \item \textbf{带中间步骤的 Transformer}：只要深度超过一个\textbf{常数}（与输入长度无关），就可以解决任何固有串行（inherently serial）问题。
    \item \textbf{直接输出答案的 Transformer}：要么需要巨大的深度，要么根本无法求解。
\end{itemize}

\begin{warningbox}{实践含义}
如果模型无法解决某个问题，一个有效策略是引导它生成更多的中间步骤。此外，可以调用外部工具来辅助中间步骤的生成——这正是 LLM Agent 框架的核心思想之一。
\end{warningbox}

\subsection{本章小结}

中间步骤思想经历了从训练到 fine-tuning 再到 prompting 的演进，但本质始终不变：让模型逐步推导而非直接跳到答案。理论上，中间步骤赋予了有限深度 Transformer 解决任意串行计算的能力。

